\section{Datasets and evaluation protocol}
\label{app:roster}
\subsection{Dataset definitions}
Table~\ref{tab:roster} describes the 22 prediction tasks. Here $d$ counts input
parameters, $N_H$ counts HF training rows, and $N_{L,\min}$ is the smallest LF
table, not the sum over fidelities. Query inputs are excluded from every
training fidelity and are subsequently partitioned into fitting and evaluation
sets. These counts describe available data, rather than identical exposure
under every training recipe. The seven PDE classes contain six elliptic,
two diffusion, three convection, two shock, one porous flow, five reaction
diffusion, and two wave datasets. ERA5 supplies the climate case and enters
dataset aggregation, but not PDE class aggregation.

\begin{table}[h!]
\centering\footnotesize\setlength{\tabcolsep}{3pt}
\caption{\textbf{Dataset definitions and available training data.} Query counts
include fitting pools and evaluation cases. Work grids are used for prediction
and scoring. Citations identify imported data, while uncited entries use our
simulation code.}
\label{tab:roster}\resizebox{\linewidth}{!}{\begin{tabular}{lrrrrrll}
\toprule
Dataset & $d$ & Levels & $N_{L,\min}$ & $N_H$ & Query & Work grid & Class \\
\midrule
Darcy & 16 & 3 & 400 & 400 & 100 & $128\times128$ & Elliptic \\
Eikonal & 3 & 2 & 400 & 400 & 100 & $64\times64$ & Elliptic \\
Helmholtz & 2 & 3 & 400 & 400 & 100 & $256\times256$ & Elliptic \\
Poisson I & 5 & 3 & 400 & 400 & 100 & $64\times64$ & Elliptic \\
Poisson II \citep{niu2024} & 5 & 5 & 64 & 64 & 512 & $128\times128$ & Elliptic \\
Pressure Poisson & 2 & 4 & 400 & 400 & 100 & $64\times64$ & Elliptic \\
Heat I & 3 & 3 & 400 & 400 & 100 & $64\times64$ & Diffusion \\
Heat II \citep{niu2024} & 3 & 5 & 1024 & 1024 & 512 & $128\times128$ & Diffusion \\
Cavity & 1 & 4 & 400 & 400 & 100 & $256\times256$ & Convection \\
Navier Stokes \citep{niu2024} & 2 & 2 & 256 & 256 & 256 & $64\times64$ & Convection \\
Rayleigh Bénard & 1 & 2 & 400 & 400 & 100 & $64\times64$ & Convection \\
Burgers & 2 & 3 & 400 & 400 & 100 & $256\times256$ & Shocks \\
Euler & 7 & 3 & 400 & 400 & 100 & $128\times128$ & Shocks \\
Porous medium & 2 & 3 & 400 & 400 & 100 & $256\times256$ & Porous flow \\
Allen Cahn & 19 & 3 & 400 & 400 & 78 & $256\times256$ & Reaction diffusion \\
Cahn Hilliard I & 2 & 2 & 400 & 400 & 100 & $64\times64$ & Reaction diffusion \\
Cahn Hilliard II & 19 & 3 & 400 & 400 & 100 & $256\times256$ & Reaction diffusion \\
Fisher KPP & 50 & 3 & 400 & 400 & 100 & $256\times256$ & Reaction diffusion \\
Phase field crystal & 18 & 3 & 400 & 400 & 100 & $128\times128$ & Reaction diffusion \\
Shallow water & 2 & 3 & 400 & 400 & 100 & $128\times128$ & Waves \\
Wave & 3 & 2 & 400 & 400 & 100 & $80\times80$ & Waves \\
ERA5 \citep{hersbach2020,niu2024} & 12 & 9 & 313 & 55 & 17 & $128\times256$ & Climate \\
\bottomrule
\end{tabular}
}
\end{table}

Poisson II, Heat II, and Navier Stokes use data released with MFRNP
\citep{niu2024}. ERA5 combines its climate model sources with reanalysis targets
\citep{hersbach2020,niu2024}. Other entries were simulated for this collection.
Pressure Poisson adapts an analytic construction motivated by \citet{partin2022},
Euler a quadrant setup motivated by The Well \citep{ohana2024}, and Shallow
water a radial dam break recipe motivated by PDEBench \citep{takamoto2022}.
These recipe influences do not make the independently simulated entries copies
of those benchmarks. Euler, Burgers, and Shallow water call PyClaw
\citep{ketcheson2012pyclaw}, using first order Godunov methods at coarse
fidelities and SharpClaw WENO reconstruction at the fine fidelity. Our dam
break varies both the inner water height and radius on a unit square, with
independent solves at each grid. Phase field crystal implements the conserved
dynamics of \citet{elder2004}. Roman numerals distinguish tasks with the same equation.
Poisson I and II have related parameterizations, while Cahn Hilliard I and II
have 2 and 19 input parameters. Cahn Hilliard II, Allen Cahn, and Fisher KPP use
$64\times64$, $128\times128$, and $256\times256$ fidelity grids. Phase field
crystal uses $32\times32$, $64\times64$, and $128\times128$ grids. Their input
dimensions are 19, 19, 50, and 18, respectively. Heat fields span space and
time, and Helmholtz predicts amplitude normalized fields.

\phantomsection
\label{app:generatorchecks}
\noindent\textbf{Numerical checks and their scope.}
Verification scripts test the custom implementations against analytical
solutions, conservation laws, and refined numerical solutions. In a one
dimensional Barenblatt test \citep{vazquez2007}, the Porous medium solver's
relative $L^2$ error falls from $8.1\times10^{-4}$ at 128 points to
$2.6\times10^{-5}$ at 512 points. An Allen Cahn mesh check reduces error
from $4.9\times10^{-5}$ at 64 points to $2.3\times10^{-6}$ at 256 points
against a 512 point reference. The archived checks also test mass conservation
for Cahn Hilliard II and Phase field crystal, and the Helmholtz discrete
equation residual. These are controlled checks of selected solver components,
including one dimensional tests for solvers used in two dimensions. They do
not certify every sampled parameter or the accuracy of every stored field.
The companion package records the scripts, measured errors, and scope of
each check. Published software and recipe references establish provenance,
while the implementation checks are our own.

Related archives can share physical simulations, so the 22 tasks are not
independent samples of PDE families. Three entries have additional limitations:
Poisson II has inherited numerical defects, Cahn Hilliard I omits realization
information from its parameters, and Pressure Poisson uses synthetic coarse
fields. Appendix~\ref{app:subsets} checks sensitivity to these entries without
certifying the other solvers.

\subsection{Training, fitting, and evaluation cases}
\label{app:splits}
Reserved fitting and evaluation inputs are excluded from LF and HF training
before normalization, basis construction, or model fitting. Within each PDE
partition, input hashes separate fitting from evaluation and define nested
budgets $K\in\{5,10,20\}$. Every rule uses the same fitting allowance, including
any selection or tuning. Models remain fixed across three partitions, which
measure sensitivity to case assignment rather than variation across training
seeds. Some cases change roles between partitions. Across the 21 PDE datasets,
3,058 query rows produce 1,842 evaluation appearances corresponding to 1,505
distinct dataset and row pairs. Shared simulations across related datasets can
further reduce physical independence. Historical evaluation results influenced
development and reporting scope, making these comparisons retrospective.

Allen Cahn, Cahn Hilliard II, Fisher KPP, and Phase field crystal each supply
400 training inputs per fidelity. Their query pools contain 78, 100, 100, and
100 inputs. Each partition reserves 16 evaluation cases for Allen Cahn and
20 for each of the other three datasets. ERA5 has 55 HF training inputs,
ten fitting pool inputs, and seven evaluation inputs, with a smallest retained
LF table of 313 rows. The additional Heat I inputs in
Appendix~\ref{app:additionalheat} form a separate supplementary evaluation.
The subset and loss diagnostics cover the original 17 PDE tasks, while the
primary comparison covers all 22 datasets.

Predictions share target grids, scaling, and input identities within each
comparison. Source hashes, saved target identities, or checkpoint evaluation
verify correspondence. Saved weights are fixed before evaluation answers are
scored, and replay checks reproduce individual errors, mixture errors, and
aggregate denominators. ERA5 answer files are separated from model training,
and changing evaluation answers leaves fitted weights unchanged in the
collector check. These checks verify the implemented separation without
reconstructing all prior researcher access. Error matrices use float64 and
the same stored model predictions are reused within each experiment.
Precision and checkpoint provenance records remain in the companion package.

\subsection{How the fields are displayed}
\label{app:gallery}
Figure~\ref{fig:datasets} shows row zero of each HF training table, chosen
independently of prediction error. Its grids are native array grids, while
the inventory gives the working grids used for scoring. Heat axes span space
and time. The cavity target is vorticity, with the top lid moving right and
array rows increasing upward. Its display uses a symmetric logarithmic scale,
linear within $[-1,1]$ and bounded by $[-400,400]$, so wall values do not hide
the interior structure. No values are clipped or modified. Enlarged panels,
numerical display scales, extraction metadata, and source hashes for all
22 datasets are included in the companion package.

Figure~\ref{fig:fieldcomparison} illustrates Heat I, selected after inspecting
dataset results. It uses the first evaluation case of the first partition,
not the case with the greatest improvement, and reuses the saved five example
fits. All predicted fields share a linear color scale. Absolute errors share
a logarithmic scale that is linear between zero and $10^{-5}$. Scores use
the complete $64\times64$ fields without smoothing. The selected model is M4,
so its column repeats that model. Figure~\ref{fig:overview}'s ERA5 display
uses a smoothed coarse thumbnail and interpolates the predicted working grid
to the native fine grid for illustration, without changing scored fields.
